\begin{abstract}
An AI agent can follow every instruction it was given and still leave its principal worse informed than the
principal assumes. We show this in a setting where the omission is unambiguous and costly: a delegated
negotiation in which the principal's own requirement rules out a better deal, and the agent decides whether to
mention it. Agents almost never mention it. Across \nModels\ models from \nFamilies\ families and \nEpisodes\
negotiations under a prespecified analysis plan, agents respected the requirement (\violationCount\ violations)
and protected a confidential figure (\leakCount\ leaks), yet conveyed the blocked better deal in \doneToldAll\
of the episodes where one existed; \doneLowerAll\ of reports refer to no other offer or option at all, a floor
that uses no model judgment. The requirement cost the principal \shadowMean\ points on average, and the reports
do not say so. The silence is not ignorance: asked directly afterwards, the same agents named the blocked offer
and its value in \probeKnewAll\ of the silent cases. Nor is it strategic---disclosure would have excused a
weaker outcome, and cheap talk gives an agent whose interests are aligned with its principal's no reason to
withhold. What changes the behaviour is the question rather than the instruction: rewording how the constraint
was described moved nothing, expecting the transcript to be read moved almost nothing, but asking the agent
what it had turned down cut non-disclosure from \intOpenJudged\ to \intTradeoffsJudged. Stating the deal's value
correctly tracked model capability; volunteering the trade-off did so only within one developer's ladder and was
flat within the other two. Oversight of delegated agents cannot assume the agent will volunteer what its
principal needs to know.
\end{abstract}
